
\subsubsection{6.1 Interpretation of Findings}

\textbf{Finding 1: HF Hub and OpenML exhibit domain-temporal selection bias, not random coverage gaps.}

The 30\% HF Hub coverage is concentrated in NLP benchmarks and small-scale image classification datasets adopted by the HF ecosystem (CIFAR-10, CIFAR-100, SVHN). The 22\% OpenML coverage is predominantly classical tabular datasets, with IMDB, 20 Newsgroups, and MNIST also present. Large-scale DL vision (ImageNet, COCO, Pascal VOC, KITTI), speech (LibriSpeech, TIMIT, Wall Street Journal), and graph (Cora, Citeseer) datasets --- which dominate Raff's post-2012 deep learning papers --- are absent from both platforms. This is not a data-quality problem addressable by better pipeline implementation; it is a structural consequence of each platform's design scope.

For researchers planning HF+OpenML-based reproducibility auditing of pre-2018 ML literature: the analysis will represent NLP, classical ML, and small-scale image papers, while systematically omitting the large-scale DL vision papers that constitute a substantial portion of the deep learning-era corpus. This bias operates silently unless coverage is explicitly measured for the target corpus, as done here.

\textbf{Finding 2: Retroactively-created HF cards are systematically thin.}

The 0.41 mean field-presence score indicates that even when a card exists, the fields most relevant to the misuse detection mechanism are absent. The \texttt{intended\textbackslash \_use} and \texttt{out\textbackslash \_of\textbackslash \_scope\textbackslash \_use} fields --- which directly encode scope boundaries --- are systematically missing from retroactively-created cards for foundational benchmarks such as MNIST, CIFAR-10, and IMDB. This is consistent with the timing of the Datasheets schema \citep{Gebru2021Datasheets}: cards created before 2021 were not required to conform to the full schema, and retroactive curation has been incomplete. The schema was designed for new dataset creation rather than retroactive annotation of pre-2021 benchmarks.

\textbf{Finding 3: The theoretical hypothesis remains scientifically plausible but untested.}

The three-step theoretical mechanism --- (1) high benchmark concentration signals artifact accumulation (D'Amour et al., 2022); (2) absent \texttt{intended\textbackslash \_use} and \texttt{out\textbackslash \_of\textbackslash \_scope\textbackslash \_use} fields indicate elevated out-of-context application risk \citep{Gebru2021Datasheets}; (3) the joint effect predicts reproducibility failure above paper-quality controls --- is grounded in prior literature but was not empirically tested. H-M1, H-M2, and H-M3 remain blocked. No claim about whether documentation completeness or benchmark concentration actually predicts reproducibility failure is supported by this study's data.

\textbf{Finding 4: Papers With Code is the most appropriate primary data source for a redesigned study.}

Papers With Code explicitly tracks paper-dataset linkages across all dataset domains and historical periods. Unlike HF Hub (NLP-centric) and OpenML (tabular-centric), Papers With Code is designed for the paper-dataset-code linkage required for this class of study and covers Raff's corpus temporal period. This is a theoretical recommendation based on the documented scope of Papers With Code; it has not been empirically verified. A redesigned H-E1 using the Papers With Code API would test whether coverage is adequate before proceeding to regression analysis.

\subsubsection{6.2 Limitations}

\textbf{Limitation 1: The original research question remains unanswered.} The central hypothesis --- do metadata-observable dataset misuse signals predict ML reproducibility failure? --- is neither confirmed nor refuted by this study. This paper reports an infrastructure characterization. Running regression on 30\% independent variable coverage would produce biased, uninterpretable results; blocking was the correct methodological decision.

\textbf{Limitation 2: H-E1 used canonical dataset names only.} Alternate name variants (e.g., ''uoft-cs/cifar10'', ''torchvision/cifar10'') may recover additional HF cards. The 30\% figure may represent a lower bound. However, the large-scale DL vision, speech, and graph gap is structural --- HF Hub's architecture and community do not serve these datasets under any naming convention --- and the canonical name interface is the appropriate choice for a scalable automated pipeline.

\textbf{Limitation 3: Results are specific to the Raff 2019 corpus.} Post-2019 papers with more NLP and HF-native datasets would show substantially higher HF Hub coverage. API coverage gaps may differ for other corpora, and the characterization presented here should not be generalized beyond the Raff 1984--2017 temporal and domain scope.

\textbf{Limitation 4: The theoretical mechanism remains empirically unverified.} The causal chain posited in Section 2 --- concentration $\rightarrow$ artifact accumulation; incompleteness $\rightarrow$ out-of-context application; joint effect $\rightarrow$ reproducibility failure --- is theoretically motivated but not confirmed by this study's experiments. Empirical verification requires the redesigned data infrastructure.

\textbf{Limitation 5: Single annotator ground truth.} Raff's reproducibility labels are from a single researcher's manual assessment. Single-annotator bias may systematically affect which paper types are labeled as reproducible or not. This limitation applies to any downstream study using Raff's labels, not specifically to the infrastructure characterization reported here.

\subsubsection{6.3 Future Work}

The highest-priority next experiment is a redesigned H-E1 using Papers With Code API as the primary data source, measuring whether PwC coverage for the Raff corpus meets the thresholds required for valid regression analysis. If PwC coverage is adequate ($\geq 70$\%), the original research question --- do metadata-observable misuse signals predict reproducibility failure? --- becomes testable.

Additional directions include: systematic exploration of HF Hub name variants to determine whether the 30\% figure is a lower bound; extension to post-2019 corpora where HF card coverage is substantially higher and multi-annotator reproducibility labels are available (ML Reproducibility Challenge); and a multi-source hybrid approach combining HF Hub, OpenML, Papers With Code, and arXiv metadata extraction.

